\section{Discussion}
\label{sec:discussion}

The analysis separates three objects that can easily be conflated: technological compatibility, government incentives, and formal coalition stability. A private workaround directly changes only the first. Its effect on the second depends on which rents survive the workaround, and its effect on the third additionally depends on the coalition institution.

\subsection{Economic interpretation}
\label{subsec:discussion-interpretation}

The residual-rent decomposition gives a compact interpretation of the canonical mechanism. Before private adoption, a regional member's advantage over international standardization is $\mathscr E-\mathscr D$. The outsider's adoption removes part of the member-market component, but with incomplete adaptation it need not remove all of it. The surviving amount is $R(d,v)$. The political effect is therefore governed by the position of $R(d,v)$ relative to two endogenous thresholds rather than by the existence of interoperability itself.

This distinction makes the sign ambiguity economically useful. If residual adaptation costs are sufficiently low, the outsider can become competitively effective in the bloc while members retain their disadvantage outside it; broader formal harmonization becomes relatively more attractive and may become strictly preferred. With more residual rent, the same private action can still reduce the value of regional exclusion without crossing the government's ranking boundary. With still more residual rent, the private response can make regional standardization relatively more attractive than before. Private adaptation is therefore neither a general political substitute for nor a general political complement to formal harmonization.

The fixed adoption cost $F$ plays a different role. It selects which private continuation state is reached. A low residual marginal cost $d$ is irrelevant for the government mechanism if the outsider never adopts, and a favorable government comparison is irrelevant if the firm game does not select outsider-only adaptation. This is why the joint open-set result imposes the adoption and government inequalities on the same primitives.

\subsection{What the robustness results imply}
\label{subsec:discussion-robustness}

The negative portability results discipline the interpretation. At $v=0$, private one-way adoption can still occur while the initial SU advantage disappears. Under the pre-specified singleton-network alternative, the same separation arises. Under the differentiated-demand test, outsider-only adoption again survives while the member-market political term changes sign in both Cournot and Bertrand versions.

These results suggest a hierarchy rather than a binary verdict on robustness. The scope-versus-fixed-cost logic behind asymmetric adoption is broader than the government-ranking reversal. The latter requires a product-market environment in which regional exclusion initially creates a sufficiently valuable member-market component. The canonical Cournot specification supplies one such environment and yields exact success and failure boundaries inside it; R5 shows that this condition cannot simply be presumed under another demand system.

The coalition analysis adds another layer. The intermediate unique-IS application is based on strict payoff differences and survives both strict and weak blocking under the certified asymmetry perturbation. By contrast, the high-$F$ multiplicity of symmetric regional unions relies on an indifference that disappears under weak blocking or small market-size asymmetry. The economically robust statement is therefore about continuation-payoff changes, not about one particular stable-set correspondence.

\subsection{Institutional interpretation and empirical use}
\label{subsec:discussion-institutions}

The formal partitions are stylized representations of fragmented, regional, and common standards regimes. Private adoption can be read as redesign, multi-standard production, dual certification, interface adaptation, or another costly action that allows a firm to support an additional formal standard. The model does not identify which real-world technology corresponds exactly to $F$ or $d$. It isolates the economic distinction between a fixed investment needed to support another standard and the residual per-unit burden that may remain after doing so.

This separation suggests measurable empirical objects. A practical application could distinguish the setup cost of adding a standard or certification from the residual variable cost of serving that market after adaptation. The theory predicts that changes in the latter matter for political incentives only through the continuation rents they leave in the protected markets. Testing that prediction would require institution-specific measures of firm adaptation, competitive effects, and government preferences; the present paper provides no empirical identification strategy.

The coalition result should likewise not be read as a welfare theorem. In the certified intermediate region, IS Pareto-dominates the other four formal partitions within the model's five-partition menu. No unrestricted planner problem is solved, and nothing here establishes that international standardization is globally first best.

\subsection{Limitations}
\label{subsec:discussion-limitations}

The model has three countries, one firm per country, segmented markets, linear demand, binary standard adoption, exogenous fixed costs, and a small formal coalition menu. The residual-cost extension keeps compatibility binary even when cost relief is incomplete. Endogenous compatibility intensity, heterogeneous firms, entry, investment dynamics, and bargaining could change both adoption incentives and government payoffs.

The differentiated-demand exercise is intentionally one portability test rather than a search over specifications until the result reappears. It therefore demonstrates non-portability to that model, not nonexistence under all differentiated-product or Bertrand environments. Similarly, the singleton-network exercise tests one alternative and does not establish a universal theorem over network formulations.

Finally, coalition stability is static. The analysis does not model accession negotiations, farsighted deviations, transfers, repeated interaction, or sunk adoption investments carried across institutional transitions. Those extensions could change the mapping from continuation payoffs to stable formal arrangements even when the product-market mechanism remains the same.
